\section*{Bab \ref{ch:g7:shadows-eclipses} --- Bayang-Bayang dan Gerhana}

\begin{solution}{exo:g7:shadows-eclipses:1}
\omterm{prop:g7:shadows-eclipses:umbra}{Umbra}: sumbernya tersembunyi seluruhnya --- \omterm{def:g1:light-and-shadows:source}{cahaya} dari tidak satu
bagian pun darinya yang tiba. \omterm{prop:g7:shadows-eclipses:umbra}{Penumbra}: sumbernya tersembunyi sebagian
--- \omterm{def:g1:light-and-shadows:source}{cahaya} dari sebagiannya masih tiba, sambil mengintip di atas tepi
penghalangnya.
\end{solution}

\begin{solution}{exo:g7:shadows-eclipses:2}
Sebuah titik tidak mempunyai ``bagian'' untuk disembunyikan sebagian:
setiap tempat entah melihatnya entah tidak --- \omterm{def:g1:light-and-shadows:shadow}{bayang-bayang} yang
tajam, tanpa pinggiran. Langit yang mendung adalah sumber yang begitu
lebar sehingga hampir setiap bintik melihat sebagian besarnya: hampir
seluruhnya \omterm{prop:g7:shadows-eclipses:umbra}{penumbra}, dan \omterm{def:g1:light-and-shadows:shadow}{bayang-bayangnya} terhapus.
\end{solution}

\begin{solution}{exo:g7:shadows-eclipses:3}
Sinar dari \emph{tepi} sumbernya, yang menyerempet sisi penghalangnya.
\omterm{prop:g7:shadows-eclipses:umbra}{Umbranya} adalah daerah yang diluputkan sinar kedua tepinya.
\end{solution}

\begin{solution}{exo:g7:shadows-eclipses:4}
Matahari: Matahari--Bulan--Bumi, dalam urutan itu --- \omterm{def:g2:moon-in-the-sky:moon}{Bulan} harus
berupa bulan baru (di antara kita dan Matahari). \omterm{def:g2:moon-in-the-sky:moon}{Bulan}:
Matahari--Bumi--Bulan --- \omterm{def:g2:moon-in-the-sky:moon}{Bulan} harus purnama (berseberangan dengan
Matahari).
\end{solution}

\begin{solution}{exo:g7:shadows-eclipses:5}
Totalitas: hanya mereka yang berada di dalam bintik kecil \omterm{prop:g7:shadows-eclipses:umbra}{umbranya}
yang melesat. Gigitan sebagian: kawasan \omterm{prop:g7:shadows-eclipses:umbra}{penumbra} yang luas di
sekelilingnya. Sama sekali tidak melihat apa-apa: setiap orang yang
diluputkan daerah \omterm{def:g1:light-and-shadows:shadow}{bayang-bayangnya} --- sebagian besar \omterm{def:g4:solar-system:planet}{planetnya}.
\end{solution}

\begin{solution}{exo:g7:shadows-eclipses:6}
\omterm{def:g7:shadows-eclipses:lunar}{Gerhana bulan} adalah \omterm{def:g2:moon-in-the-sky:moon}{Bulan} itu sendiri yang meredup --- sebuah
perubahan pada sebuah benda yang dapat dilihat seluruh sisi malamnya
sekaligus. Totalitas matahari menuntut berdiri \emph{di dalam} \omterm{prop:g7:shadows-eclipses:umbra}{umbra}
\omterm{def:g2:moon-in-the-sky:moon}{Bulan} yang ramping di tempat ia menyentuh tanah: sebuah bintik, bukan
peristiwa selebar langit.
\end{solution}

\begin{solution}{exo:g7:shadows-eclipses:7}
\omterm{def:g6:sun-earth-moon:axisorbit}{Orbit} \omterm{def:g2:moon-in-the-sky:moon}{Bulan} miring kira-kira lima derajat terhadap bidang
Bumi--Matahari: kebanyakan bulan baru dan \omterm{ex:g2:moon-in-the-sky:shapes}{bulan purnama} lewat di atas
atau di bawah garisnya yang persis, dan \omterm{def:g1:light-and-shadows:shadow}{bayang-bayangnya} luput. Hanya
ketika fasenya dan penyeberangan bidangnya berbarengan, ketiga
dunianya berbaris.
\end{solution}

\begin{solution}{exo:g7:shadows-eclipses:8}
Proyeksi lubang jarum ke atas kertas, dan \omterm{ex:g2:moon-in-the-sky:shapes}{sabit} di tanah milik pohon
yang rimbun (atau sebuah saringan); kacamata gerhana bersertifikat
yang tidak rusak untuk memandang langsung. Mata telanjang: hanya
selama totalitasnya sendiri, dan serpihan pertama yang kembali
mengakhiri izinnya.
\end{solution}

\begin{solution}{exo:g7:shadows-eclipses:9}
Matahari kira-kira $400$ kali lebar \omterm{def:g2:moon-in-the-sky:moon}{Bulan} pada kira-kira $400$ kali
jaraknya, sehingga kedua cakramnya berpadanan di langit dengan hampir
tepat. Anak kepadanannya: totalitas memang ada tetapi hanya nyaris
saja --- singkat, dengan bintik \omterm{prop:g7:shadows-eclipses:umbra}{umbra} yang mungil; dan ketika \omterm{def:g2:moon-in-the-sky:moon}{Bulan}
menunggangi bagian jauh \omterm{def:g6:sun-earth-moon:axisorbit}{orbit} lonjongnya, cakramnya yang sedikit lebih
kecil meninggalkan lingkar tepi Matahari terbuka --- itulah cincin
gerhana cincinnya.
\end{solution}

\begin{solution}{exo:g7:shadows-eclipses:10}
Keliman \omterm{prop:g7:shadows-eclipses:umbra}{penumbra} setiap \omterm{def:g1:light-and-shadows:shadow}{bayang-bayang} digambar oleh tepi Matahari, dan
pembangunannya menunjukkan kelimannya melebar seiring selang di antara
penghalang dan layarnya: sinar dari tepi matahari yang berseberangan
merenggang makin jauh makin panjang perjalanannya melewati
penghalangnya. Kakinya berdiri pada selang nol --- tajam;
\omterm{def:g1:light-and-shadows:shadow}{bayang-bayang} kepalanya terletak dua \omterm{def:g2:measuring-length:units}{meter} perenggangan jauhnya ---
lembut.
\end{solution}

\begin{solution}{exo:g7:shadows-eclipses:11}
Tiap celah daunnya adalah kamera lubang jarum alami; tiap bercak
tanahnya adalah \omterm{prop:g5:mirrors-reflection:image}{bayangan} Matahari yang dibuat lubang jarum itu ---
bundar pada hari biasa. Selama gerhana sebagiannya Matahari sendiri
berupa \omterm{ex:g2:moon-in-the-sky:shapes}{sabit}, sehingga tiap kamera kecil yang setia memproyeksikan
sebuah \omterm{ex:g2:moon-in-the-sky:shapes}{sabit} --- dan semuanya menunjuk ke arah yang sama karena
semuanya menyalin satu \omterm{ex:g2:moon-in-the-sky:shapes}{sabit} yang sama di langit.
\end{solution}

\begin{solution}{exo:g7:shadows-eclipses:12}
\omterm{def:g1:light-and-shadows:source}{Cahaya} matahari yang menyerempet lingkar tepi Bumi menyeberangi
\omterm{def:g5:air-pressure-weather:pressure}{atmosfernya} di sepanjang jalan terpanjang yang mungkin --- jalan
sepanjang \omterm{def:g1:day-and-night:daynight}{matahari terbenam} di sekeliling \omterm{def:g4:solar-system:planet}{planetnya} sekaligus.
\omterm{def:g2:air-around-us:air}{Udaranya} membaurkan sebagian besar \omterm{def:g1:light-and-shadows:source}{cahaya} yang lewat itu (jatah yang
melukis langit Bumi), dan \omterm{def:g7:light-sources-propagation:diffusion}{pembauran} itu merampok warna berjalan pendek
paling dahulu, sehingga \omterm{def:g1:light-and-shadows:source}{cahaya} yang selamat melewati perjalanan
menyerempet yang panjang itu lalu membelok terus ke \omterm{def:g2:moon-in-the-sky:moon}{Bulan} adalah yang
selamat dari \omterm{def:g1:day-and-night:daynight}{matahari terbenam}: merah. \omterm{def:g2:moon-in-the-sky:moon}{Bulan} yang tergerhana bernyala
oleh \omterm{def:g1:light-and-shadows:source}{cahaya} setiap \omterm{def:g1:day-and-night:daynight}{matahari terbit} dan terbenam di Bumi yang jatuh
padanya bersama-sama.
\end{solution}

\begin{solution}{pb:g7:shadows-eclipses:1}
\textbf{1.} Jalurnya adalah \omterm{def:g6:describing-motion:trajectory}{lintasan} \omterm{prop:g7:shadows-eclipses:umbra}{umbra} \omterm{def:g2:moon-in-the-sky:moon}{Bulan} yang melesat
menyeberangi tanahnya; kawasan lebar yang mengapitnya adalah jejak
kaki \omterm{prop:g7:shadows-eclipses:umbra}{penumbranya}.
\textbf{2.} Gerhana sebagian: kotanya duduk di dalam \omterm{prop:g7:shadows-eclipses:umbra}{penumbranya} ---
sebuah gigitan dari Matahari, yang ditonton lewat pelindung yang
sepatutnya, dan tanpa kegelapan.
\textbf{3.} Bahkan serpihan Matahari lima persen pun ribuan kali lebih
terang daripada \omterm{def:g1:light-and-shadows:source}{cahaya} \omterm{ex:g2:moon-in-the-sky:shapes}{bulan purnama}: ia menjaga langitnya tetap
tersinari dan koronanya tetap tak terlihat. Malam-pada-tengah-hari
menjadi milik \omterm{prop:g7:shadows-eclipses:umbra}{umbranya} saja --- tidak ada ``hampir totalitas''.
\textbf{4.} Bintik \omterm{prop:g7:shadows-eclipses:umbra}{umbranya} kasarnya bundar: sebuah perkemahan pada
garis tengahnya duduk di bawah lebar penuh bintiknya selagi ia melesat
lewat, sedangkan perkemahan dekat tepi jalurnya hanya menangkap
setali busur lingkarannya --- penyeberangan yang lebih pendek,
totalitas yang lebih pendek.
\textbf{5.} Kacamata gerhana bersertifikat menghalangi bakaran
Mataharinya; kacamata hitam tidak --- ia meredupkan silaunya sambil
meloloskan jatah yang membakar selaput jala mata, sehingga menatap
dengan nyaman menjadi mungkin: lebih buruk daripada tidak ada.
\textbf{6.} Lubang jarum di satu ujung, layar kertas kalkir di ujung
yang lain (atau kertas polos di tanah): berdirilah dengan
\emph{punggungmu} menghadap Matahari, lubangnya dibidikkan lewat
bahumu, lalu tontonlah cakram yang terproyeksi --- jangan pernah lewat
lubangnya.
\textbf{7.} Ditolak: pantulan airnya adalah \omterm{prop:g5:mirrors-reflection:image}{bayangan} \omterm{def:g5:mirrors-reflection:reflection}{cermin} Matahari,
hampir sama menyilaukannya dengan Matahari sendiri --- bakaran yang
dipantulkan tetap saja bakaran.
\textbf{8.} Berpuluh-puluh \omterm{prop:g5:mirrors-reflection:image}{bayangan} \omterm{ex:g2:moon-in-the-sky:shapes}{sabit} pada seprainya: tiap
lubangnya adalah kamera lubang jarum, dan masing-masing memproyeksikan
Matahari yang tergigit --- sebuah saringan penuh gerhana.
\textbf{9.} \omterm{def:g2:moon-in-the-sky:moon}{Bulan} --- tidak terlihat karena silaunya sampai tepi
cakramnya mulai menutupi cakram Matahari; perkemahannya telah masuk ke
\omterm{prop:g7:shadows-eclipses:umbra}{penumbranya}.
\textbf{10.} Matahari yang terlihat sedang menyusut menuju serpihan
yang tipis --- sumber yang kian serupa titik. Sumber lebih kecil,
keliman \omterm{prop:g7:shadows-eclipses:umbra}{penumbra} lebih sempit: tepi lembut setiap \omterm{def:g1:light-and-shadows:shadow}{bayang-bayang}
mengetat, dan dunianya berubah aneh dan tajam.
\textbf{11.} \omterm{prop:g7:shadows-eclipses:umbra}{Umbra} \omterm{def:g2:moon-in-the-sky:moon}{Bulan} itu sendiri --- kelas itu berdiri di dalam
ujung kerucut \omterm{def:g1:light-and-shadows:shadow}{bayang-bayangnya} selagi ia melesat di atas perkemahannya.
Mengelilingi cakram hitamnya tergantung koronanya, \omterm{def:g5:air-pressure-weather:pressure}{atmosfer} luar
Matahari yang berkilau mutiara, yang terlihat hanya sekarang;
serpihan pertama cakram yang kembali memulihkan aturan \omterm{def:g1:light-and-shadows:source}{cahaya} siang
sepenuhnya --- kacamatanya dipasang lagi.
\textbf{12.} Misalnya: ``Gerhana adalah \omterm{def:g1:light-and-shadows:shadow}{bayang-bayang} dengan dunia
sebagai lampu, penghalang, dan layarnya: \omterm{prop:g7:shadows-eclipses:umbra}{umbra} \omterm{def:g2:moon-in-the-sky:moon}{Bulan} yang menyapu
tanah kita, atau \omterm{def:g2:moon-in-the-sky:moon}{Bulan} yang berlayar menembus \omterm{prop:g7:shadows-eclipses:umbra}{umbra} kita --- geometri
yang dapat kamu gambar dengan dua sinar dan sebuah penggaris, yang
dipentaskan pada ukuran termegah di alam.''
\end{solution}
